% Authoritative exercise source. Preserve all derivations and diagrams.

\ExerciseStart{G3.10}{Usare e ripulire un'ancilla}{Consolidamento}

\BookHeading{Problema, dati e obiettivo}

Preparare \(\ket{++}\ket0\), calcolare AND in un’ancilla, applicarle Z e ripristinarla con la stessa Toffoli. Confrontare la successiva interferenza dei primi due qubit con il caso in cui il ripristino sia omesso.

\ExTheory{\ExRef{c3-ausiliari}{3.09}}{La Toffoli scrive \(ab\) nell’ausiliario inizialmente zero; Z introduce un segno sul valore uno. Ripetere la Toffoli annulla il contenuto calcolato, perché \(ab\oplus ab=0\), lasciando il segno sul ramo pertinente.}

\ExPrediction{Dopo il ripristino tutti i rami dovrebbero condividere la stessa ancilla. Se la traccia del ramo 11 rimanesse nell’ausiliario, la successiva interferenza del registro dati dovrebbe cambiare.}

\ExFigure{\begin{tikzpicture}[>=Stealth,exlabel]
\node[exbox] (n0) at (0.0,0.0) {calcola AND};\node[exbox] (n1) at (3.6,0.0) {applica $Z$};\draw[exother] (n0)--(n1);\node[exbox] (n2) at (7.2,0.0) {ripristina AND};\draw[exother] (n1)--(n2);
\end{tikzpicture}}{La marcatura coerente richiede il calcolo e la sua inversa.}


\begin{ExTasks}

\item Sviluppare \(\ket{++}\ket0\).

\item Calcolare la Toffoli sui quattro rami.

\item Applicare Z all’ancilla.

\item Ripristinare e fattorizzare l’ausiliario.

\item Applicare H ai due dati e calcolare \(P(00)\) nei due casi.

\end{ExTasks}

\ExSolution

\ExStep{Preparare tutti i rami}

\ExEq{\ket{\Psi_0}&=\tfrac12(\ket{000}+\ket{010}+\ket{100}+\ket{110}).}
Ogni ramo ha probabilità \(1/4\).

\ExStep{Calcolare AND}

\ExEq{\ket{\Psi_1}&=\tfrac12(\ket{000}+\ket{010}+\ket{100}+\ket{111}).}
Soltanto il ramo con dati 11 cambia l’ausiliario.

\ExStep{Marcare con una fase}

\ExEq{\ket{\Psi_2}&=\tfrac12(\ket{000}+\ket{010}+\ket{100}-\ket{111}).}
Z lascia i rami con ancilla zero invariati e moltiplica per \(-1\) quello con ancilla uno.

\ExStep{Ripristinare l’ausiliario}

\ExEq{\ket{\Psi_3}&=\tfrac12(\ket{000}+\ket{010}+\ket{100}-\ket{110})\\
&=\tfrac12(\ket{00}+\ket{01}+\ket{10}-\ket{11})\otimes\ket0.}
Il segno resta nel registro dei dati.

\ExStep{Interferire dopo il ripristino}

Applicando \(H\otimes H\) al vettore \((1,1,1,-1)^T/2\), si ottiene
\ExEq{\tfrac14\colv{1+1+1-1\\1-1+1+1\\1+1-1+1\\1-1-1-1}
&=\tfrac12\colv{1\\1\\1\\-1}.}
Le quattro probabilità sono \(1/4\).

\ExStep{Interferire conservando la traccia}

Se si parte da \(\Psi_2\), per l’esito dati 00 l’ampiezza con ancilla zero è \((1+1+1)/4=3/4\), mentre quella con ancilla uno è \(-1/4\). I due stati dell’ancilla sono ortogonali:
\ExEq{P(00)&=|3/4|^2+|-1/4|^2=10/16=5/8.}
Per ogni altro esito la somma sui tre rami ha modulo \(1/4\), e l’ampiezza dell’ancilla uno ha modulo \(1/4\): \(P=2/16=1/8\). Il totale è \(5/8+3/8=1\).

\par\noindent\begin{minipage}{\linewidth}
\ExFigure{\begin{tikzpicture}[>=Stealth,exlabel]
\draw[exaxis] (-.6,0)--(4.4,0) node[right] {esito};\draw[exaxis] (-.6,0)--(-.6,3.6) node[above] {$p$};\node[left] at (-.6,3) {$ 1 $};\draw[Ink!15,dashed] (-.6,3)--(3.9000000000000004,3);\filldraw[fill=Accent!25,draw=Accent] (0.0,0) rectangle (0.65,1.875);\node[below] at (0.325,0) {$ 00 $};\node[above] at (0.325,1.875) {$ 0.625 $};\filldraw[fill=Accent!25,draw=Accent] (1.1,0) rectangle (1.75,0.375);\node[below] at (1.425,0) {$ 01 $};\node[above] at (1.425,0.375) {$ 0.125 $};\filldraw[fill=Accent!25,draw=Accent] (2.2,0) rectangle (2.85,0.375);\node[below] at (2.5250000000000004,0) {$ 10 $};\node[above] at (2.5250000000000004,0.375) {$ 0.125 $};\filldraw[fill=Accent!25,draw=Accent] (3.3000000000000003,0) rectangle (3.95,0.375);\node[below] at (3.6250000000000004,0) {$ 11 $};\node[above] at (3.6250000000000004,0.375) {$ 0.125 $};
\end{tikzpicture}}{Distribuzione dopo le Hadamard quando il ripristino è omesso; con ripristino tutti gli esiti valgono \(1/4\).}


\ExInterpret{Le probabilità nella base computazionale prima della ricombinazione possono nascondere il problema. Il confronto dopo H mostra l’effetto dell’informazione residua nell’ancilla.}

\ExLab{Nella sezione 3.09 confrontare «Marcatura con ripristino» e «Informazione residua senza ripristino». Verificare l’ausiliario in ogni ramo. Le Hadamard finali dell’esercizio forniscono un controllo ulteriore, svolto qui esplicitamente.}


\end{minipage}
